\section{System Architecture}

The system has ten containers on one Docker bridge network called
\texttt{mic-net}. Containers find each other by service name through Docker's
internal DNS, so \texttt{gen-http2} can simply call
\texttt{http://target-http2:8080} without anyone configuring an IP address
anywhere.

The ten containers fall into four groups. The \textbf{controller} (FastAPI,
port 8000) is the brain: it reads the YAML profile, drives the four generators
through their REST APIs and runs the Adaptive Control loop. The four
\textbf{generators} (\texttt{gen-http2}, \texttt{gen-quic}, \texttt{gen-mqtt},
\texttt{gen-tcpudp}) each keep their own state and expose the same four
endpoints: start, stop, patch the config, get the status. They are built on
\texttt{httpx} (HTTP/2), \texttt{aioquic} \cite{aioquic} (QUIC), \texttt{paho-mqtt}
\cite{pahomqtt} (MQTT) and plain Python sockets (raw TCP/UDP). The three
\textbf{target services} (\texttt{target-http2}, \texttt{target-quic},
\texttt{target-tcpudp}) plus the Mosquitto broker \cite{mosquitto} simply receive traffic and,
in two of the three cases, report what they received back to the metrics
collector. \texttt{target-http2} runs on Hypercorn \cite{hypercorn}, the only
ASGI server in the project that supports HTTP/2 natively. The \textbf{metrics collector} (Flask, port 9090) aggregates
everything into one \texttt{/metrics} response, and the \textbf{dashboard}
(a single static HTML file served over port 3000) polls the controller every
two seconds and renders the live state.

Fig.~\ref{fig:architecture} shows how a request actually travels through the
system. The dashboard never talks to a generator directly; every change goes
through the controller, which is what makes it possible to log every
configuration change with a timestamp, as the assignment requires.

\begin{figure*}[t]
\centering
\begin{tikzpicture}[
  box/.style={draw, rounded corners, minimum width=2.6cm, minimum height=0.7cm, align=center, font=\footnotesize},
  arr/.style={-{Latex[length=2mm]}, thick},
  column sep/.style={column sep=1.1cm},
]
\node[box] (dash) at (5.2,6) {Dashboard\\:3000};
\node[box] (ctrl) at (5.2,4.6) {Controller\\:8000};

\node[box] (h2)     at (0,3) {gen-http2};
\node[box] (quic)   at (3.4,3) {gen-quic};
\node[box] (mqtt)   at (6.8,3) {gen-mqtt};
\node[box] (tcpudp) at (10.2,3) {gen-tcpudp};

\node[box] (th2)    at (0,1.4) {target-http2};
\node[box] (tquic)  at (3.4,1.4) {target-quic};
\node[box] (broker) at (6.8,1.4) {mosquitto};
\node[box] (ttcp)   at (10.2,1.4) {target-tcpudp};

\node[box, minimum width=4.2cm] (metrics) at (1.7,0) {Metrics Collector\\:9090};

\draw[arr] (dash) -- (ctrl) node[midway, right, font=\scriptsize] {GET /status};
\draw[arr] (ctrl) -- (h2);
\draw[arr] (ctrl) -- (quic);
\draw[arr] (ctrl) -- (mqtt);
\draw[arr] (ctrl) -- (tcpudp);
\draw[arr] (h2) -- (th2);
\draw[arr] (quic) -- (tquic);
\draw[arr] (mqtt) -- (broker);
\draw[arr] (tcpudp) -- (ttcp);
\draw[arr] (th2) -- (metrics);
\draw[arr] (tquic) -- (metrics);
\draw[arr, dashed] (ctrl.east) -- ++(2.2,0) |- node[pos=0.97, above, font=\scriptsize] {GET /metrics} (metrics.east);
\end{tikzpicture}
\caption{Component diagram of the single-machine deployment. Solid arrows are
REST calls or raw protocol traffic; the dashed arrow is the controller reading
aggregated metrics for status reporting and Adaptive Control.}
\label{fig:architecture}
\end{figure*}

One detail worth pointing out: \texttt{target-tcpudp} does not report to the
metrics collector at all, unlike the other two targets. It only accepts
connections and discards the data. This is deliberate, since the point of that
container is to make a TCP RST visible in Wireshark when it goes down, and
adding a reporting path would only get in the way of that.
